%% ============================================================
%% sections/s360k.tex
%% United States Code, Title 21 - § 360k. State and local requirements respecting devices
%% (FD&C Act § 521). Source identifier: /us/usc/t21/s360k
%% Relevance to the Physical AI oncology trial context:
%% State and local requirements respecting devices; device preemption, cited to keep § 360(k) and § 360k distinct.
%% Reproduced faithfully from the OLRC USLM XML (through Pub. L. 119-93).
%% No bracketed instructions or file names are inserted into the law.
%% ============================================================
\uscsection{§~360k.}{State and local requirements respecting devices}
\uscprov{1}{\textbf{(a)}\textbf{ General rule} Except as provided in subsection (b), no State or political subdivision of a State may establish or continue in effect with respect to a device intended for human use any requirement-}
\uscprov{2}{\textbf{(1)} which is different from, or in addition to, any requirement applicable under this chapter to the device, and}
\uscprov{2}{\textbf{(2)} which relates to the safety or effectiveness of the device or to any other matter included in a requirement applicable to the device under this chapter.}
\uscprov{1}{\textbf{(b)}\textbf{ Exempt requirements} Upon application of a State or a political subdivision thereof, the Secretary may, by regulation promulgated after notice and opportunity for an oral hearing, exempt from subsection (a), under such conditions as may be prescribed in such regulation, a requirement of such State or political subdivision applicable to a device intended for human use if-}
\uscprov{2}{\textbf{(1)} the requirement is more stringent than a requirement under this chapter which would be applicable to the device if an exemption were not in effect under this subsection; or}
\uscprov{2}{\textbf{(2)} the requirement-}
\uscprov{3}{\textbf{(A)} is required by compelling local conditions, and}
\uscprov{3}{\textbf{(B)} compliance with the requirement would not cause the device to be in violation of any applicable requirement under this chapter.}
\uscsource{(June 25, 1938, ch. 675, §~521, as added Pub. L. 94-295, §~2, May 28, 1976, 90 Stat. 574.)}

%% - DRAFTING INSTRUCTIONS (added for VVUQ-04 draft-bill; the reproduced statutory text above is unchanged) -
\begin{draftbox}{§ 360k (21 U.S.C. § 360k) - state and local requirements; the preemption boundary and savings clause}
\dstep{Role in the amendment: § 360k is device preemption: a State may not impose a requirement different from, or in addition to, a federal device requirement. The conforming amendment adds a savings clause so the mass-adopted State human-in-the-loop and audit pattern is preserved to the extent it is not different from, or in addition to, a federal requirement, marking the federal and State boundary without expanding preemption.}
\dstep{Critical naming: this is § 360k (the section on State and local requirements), which is distinct from § 360(k) (the 510(k) premarket notification handled in sections/s360.tex). Keep the two separate per papers/VVUQ-04/instruct-bill/10-legal-crosswalk-and-bill-style.md (Section E, item 3).}
\dstep{Amendatory action: Section 521 of the Federal Food, Drug, and Cosmetic Act (21 U.S.C. § 360k) is amended by adding at the end the following new subsection: "(c) Rule of construction for Physical AI oversight. Nothing in this section preempts a State requirement of general applicability that a licensed human retain the authority to monitor, intervene, override, or terminate an artificial intelligence or robotic system, or that an audit record be maintained, to the extent that the requirement is not different from, or in addition to, a requirement applicable under this chapter to a device, including the verification requirements of section 515D." Draft it as a savings clause only; do not enlarge or narrow the preemption rule in subsection (a).}
\dstep{Source synthesis: process papers/VVUQ-04/instruct-bill/06-state-medical-ai-laws.md (the California Physicians Make Decisions Act SB 1120; the Illinois Wellness and Oversight for Psychological Resources Act HB 1806; Texas HB 149; Maryland HB 820, as mass-adopted human-over-AI practice) and 10-legal-crosswalk-and-bill-style.md (state laws cited as findings, not as operative authority) together with papers/VVUQ-03/final-paper/sections/prior\_law.tex (the State law crosswalk). Resolve to references.bib key usc-fdca; cite the specific State statutes only in the SEC. 2 findings (via the instruct-bill state-laws.bib keys such as ca-sb1120 and il-hb1806), not in this operative subsection.}
\dstep{Comparative print rendering: reproduce subsections (a) and (b) unchanged and show the new subsection (c) wrapped in \cprinsert{...} at the end.}
\dstep{Formatting: single hyphens only; § for every codified reference; even RaggedRight spacing; reword so the savings clause does not end in a one- or two-word line.}
\end{draftbox}
